\chapter{Varying a region before and after opening a face}\label{ch:04}
Cut a vibrating string at a point and hold the displacement there equal to a
specified function of time. Each half now needs a force at the cut to follow
that history. When the two halves are rejoined, it is the \emph{sum} of the
required forces that must vanish. Setting each force to zero would describe
two free ends instead. We will obtain the corresponding statement for a field
from the first variation of its action. The calculation also tells us which
variables play the role of the displacement and where the forces are measured.

\section{The variation at prescribed data}
Let $R$ be an oriented Lorentzian slab of signature $(-,+,\ldots,+)$.
Use the same future slice orientation on the two spacelike ends, so that the
oriented boundary is $\Sigma_f-\Sigma_i+B+\Gamma$. The timelike face $B$
is a physical wall, possibly a regulator for a boundary at infinity; its
boundary condition and action are part of the theory. The timelike face
$\Gamma$ is an artificial cut. For simplicity the side faces meet the end
slices but not one another. Write $\partial_F\Sigma=\Sigma\cap F$ with the
orientation inherited as a boundary of $\Sigma$. We use $d$ for the spacetime
differential and $\delta$ for variation of fields at fixed spacetime point;
the latter commutes with spacetime differentiation in our formulas.

\begin{definition}
A regional problem at prescribed cut datum $j$ consists of a configuration
space with the physical condition on $B$ and a specified trace condition
$b_\Gamma(\phi)=j$ on $\Gamma$, together with a differentiable action on
that space. A trace assigns to a bulk field its selected boundary values;
these can include derivatives. A variation in this space satisfies
$\delta b_\Gamma=0$, not necessarily $\delta\phi|_\Gamma=0$.
\end{definition}

Dirichlet data fix the boundary field itself. Fixing a normal momentum
instead usually requires a different face term. To see why that term belongs
in the variation, write
\[
 S_R[j]=\int_R L+\int_B\ell_B+\int_\Gamma\ell_\Gamma[j],
 \qquad \delta L={\cal E}_a\,\delta\phi^a+d\Theta.
\]
Here $L$ is a spacetime top form, ${\cal E}_a$ are the Euler--Lagrange
densities and $\Theta$ is a spacetime form of one lower degree, linear in a
field variation. On each side face $F$ differentiability requires the
remainder $\Theta|_F+\delta\ell_F$ to vanish or to be a tangential derivative
$d_F C_F$ on the allowed variations. The corner potential $C_F$ has
spacetime degree $\dim R-2$ and is linear in the variation. If the wall
supports its own equation,
its coefficient is kept in the variation until that equation is imposed.
At fixed $j$, imposing the boundary conditions but not the bulk equations,
Stokes' theorem gives
\begin{align}
 \delta S_R[j]&=\int_R {\cal E}_a\,\delta\phi^a
              +\theta_{\Sigma_f}-\theta_{\Sigma_i},\nonumber\\
 \theta_\Sigma&=\int_\Sigma\Theta
         -\sum_{F=B,\Gamma}\int_{\partial_F\Sigma} C_F.
 \label{eq:04-endpoint}
\end{align}
The minus sign can be found without a metric: at the final corner,
$\partial F$ induces the opposite orientation to $\partial_F\Sigma_f$,
since the boundary of the oriented boundary vanishes. A tangential derivative
on a wall therefore contributes to the endpoint potential; it is not zero.

For a concrete corner, take a flat strip with a real scalar and add a kinetic
term $\tfrac12\kappa\int_{t_i}^{t_f}\dot q_B^2\,dt$ on one physical end of
the string, where $q_B=\phi|_B$. Its variation is
\[
 -\kappa\int_{t_i}^{t_f}\ddot q_B\,\delta q_B\,dt
      +[\kappa\dot q_B\,\delta q_B]_{t_i}^{t_f}.
\]
For the bulk kinetic term used below, the wall equation with free $q_B$ is
$-\partial_n\phi-\kappa\ddot q_B=0$, where $n$ points out of the strip.
Even on a solution of this equation, the endpoint term is still present: the
wall carries momentum $\kappa\dot q_B$. It must not be discarded with the
wall equation.

\section{Releasing a common trace}
To open $\Gamma$, allow $j$ to range over the admissible histories and keep
the condition on $B$ fixed. If $\mathcal C_R[j]$ is the configuration space
at prescribed $j$, this gives $\mathcal C_R^{\mathrm{open}}
=\bigcup_{j\;\mathrm{admissible}}\mathcal C_R[j]$; a history must still
respect any compatibility condition where the cut meets an end slice.
Remove the face action used \emph{only} to impose
the auxiliary closure, or explicitly undo its boundary Legendre transform.
Do not remove a term that defines the physical wall. Opening is not a
stationarity condition on one side of the cut. It makes the cut part of the
first variation visible as a linear functional of varying trace data $q$:
\[
 \delta S_R^{\mathrm{open}}=\int_R {\cal E}_a\delta\phi^a
 +\theta^{\mathrm{open}}_{\Sigma_f}-\theta^{\mathrm{open}}_{\Sigma_i}
 +\langle\Pi,\delta q\rangle_\Gamma.
\]
Here tangential derivatives of $\delta q$ have been integrated by parts;
the resulting corner contributions belong to $\theta^{\mathrm{open}}$.
For example, subtracting $\langle\Pi,q\rangle_\Gamma$ would turn the
face variation $\langle\Pi,\delta q\rangle_\Gamma$ into
$-\langle\delta\Pi,q\rangle_\Gamma$, suitable for fixing $j=\Pi$. Undoing
this subtraction restores the original face one-form; it does not set the
momentum to zero.
For a first-order scalar Lagrangian and Dirichlet closure, $q=j=\phi|_\Gamma$
and no artificial face action is needed. With fixed metric and volume form
$\epsilon$, the action density and its variation are
\begin{align*}
 L&=-\tfrac12(\nabla_\mu\phi\nabla^\mu\phi+m^2\phi^2)\epsilon,\\
 \delta L&=(\nabla_\mu\nabla^\mu\phi-m^2\phi)\delta\phi\,\epsilon
       -\nabla_\mu(\delta\phi\nabla^\mu\phi)\epsilon.
\end{align*}
On a timelike face the last term gives
$\langle\Pi,v\rangle_\Gamma=-\int_\Gamma (n^\mu\nabla_\mu\phi)v\,d\mu_\Gamma$,
where $v=\delta q$, $n$ is unit and outward, and $d\mu_\Gamma$ is the positive
induced measure. Thus $\Pi$ is a dual trace \emph{density} with the outward
sign already included. In arbitrary face coordinates $y$ it contains
$\sqrt{|\det\gamma|}\,|d^{\dim\Gamma}y|$, not just a component of the
normal derivative. Equivalently one can integrate the form
$-(\partial_n\phi)v\,\iota_n\epsilon$ with the outward face orientation.
Do not insert a second orientation sign when using the density formula.
One can even avoid normalizing a metric normal. For an outward conormal
$\nu$ and any transverse $N$ with $\nu(N)=1$, let
$\epsilon_\nu=(\iota_N\epsilon)|_\Gamma$. The outward surface density is
$\nu_\mu\,|\epsilon_\nu|$. Rescaling $\nu$ by a positive
function rescales $N$ inversely, and adding a tangential vector to $N$ has
zero effect on the pullback. This explains why the pairing, rather than a
bare coordinate component, is the invariant object.

Suppose a fixed smooth identification $D$ carries trace data on side 2 to
side 1 with smooth inverse. Continuity $q_1=D(q_2)$ is a restriction on
configurations, made before varying them. Differentiating $q_2=D^{-1}(q_1)$
gives $\delta q_2=T_{q_1}D^{-1}\,\delta q_1$; only for a linear map may one
use $D^{-1}$ itself on variations. At the matched data the dual transport is
defined by $\langle D^\vee\Pi_2,v\rangle_{\Gamma_1}
=\langle\Pi_2,T_{q_1}D^{-1}\,v\rangle_{\Gamma_2}$, including the change of
measure. Consequently, the sum of the two opened actions has the cut variation
\[
 \langle\Pi_1,\delta q_1\rangle_{\Gamma_1}
 +\langle\Pi_2,T_{q_1}D^{-1}\,\delta q_1\rangle_{\Gamma_2}
 =\langle\Pi_1+D^\vee\Pi_2,\delta q_1\rangle_{\Gamma_1}.
\]
Stationarity gives $\Pi_1+D^\vee\Pi_2=0$ if all compactly supported common
trace variations are available. With a restricted trace space it only says
that the sum annihilates that space. At the flat cut $x=0$, with region 1 at
$x<0$ and region 2 at $x>0$, the sum per unit measure is
$-\partial_x\phi_1+\partial_x\phi_2$: opposite outward normals give
continuity of the derivative in a \emph{single} signed coordinate.

\section{Which traces and which corner representative?}
The first variation, not a count of field components, decides how many
response densities are present. Consider a term proportional to
$\tfrac12(\partial_x^2\phi)^2$ in a flat collar, suppressing spectator
coordinates. At fixed tangential coordinates, two integrations by parts give
\[
 \int_a^b \phi''\,\delta\phi''\,dx
 =\int_a^b \phi''''\,\delta\phi\,dx
    +[\phi''\,\delta\phi'-\phi'''\,\delta\phi]_a^b.
\]
Now $q=(\phi,\partial_x\phi)$, and both coefficients must be paired in the
common signed $x$ coordinate. A smooth bulk field restricted to the two
regions matches \emph{all} normal derivatives in that coordinate. Conversely,
conditions from the displayed first variation alone guarantee neither
infinite-order matching nor the regularity of a weak solution. Such a
conclusion needs the equations and a regularity argument.

There is also a genuine freedom of presentation \cite{HW}. For a local
spacetime form $Y$ linear in the variation, $\Theta\mapsto\Theta+dY$ and
$C_F\mapsto C_F+Y|_F$ leave \eqref{eq:04-endpoint} unchanged by Stokes'
theorem. Moving a total derivative $d\mu$ into $L$ and subtracting $\mu|_F$
from each side action shifts $\Theta$ by $\delta\mu$; it changes
$\theta_\Sigma$ only by $\delta\int_\Sigma\mu$. Its field-space exterior
derivative is unchanged. An unmatched wall interaction is not such an
ambiguity: it changes what the common trace must do. The following test
detects the difference.

\begin{exercise}
On the scalar cut at $x=0$ leave the single surface action
$-\tfrac12\lambda\int j(t)^2\,dt$ in the joined system. Does stationarity
still give continuity of the normal derivative in the common coordinate?
\end{exercise}
\noindent\emph{Answer.} Its variation is $-\lambda\int j\,\delta j\,dt$.
Adding the two outward responses gives
$\partial_x\phi_2-\partial_x\phi_1=\lambda j$. The result describes a
physical defect in the string, not an invisible cut.
